Use slicing as in \cite{r-61}, Theorem 1.2(a)(i). Write \(G_\rho(\omega)=G(\rho\omega,\rho)\). The surface measure in these coordinates is \(\sqrt2\,\rho^{n-1}d\rho\,d\sigma(\omega)\). Thus
\[
 E_\Sigma G(x,t)=\sqrt2\int_1^2 e^{it\rho}\rho^{n-1}E(G_\rho)(\rho x)\,d\rho .
\]
Apply Hausdorff--Young in \(t\) and Minkowski in \(x,\rho\), using \(b\ge b'\). After spatial dilation \(\rho x\mapsto x\) this gives
\[
 \|E_\Sigma G\|_{L^b_{x,t}}\lesssim
 \left(\int_1^2\big[\rho^{n-1-n/b}\|E(G_\rho)\|_{L^b_x}\big]^{b'}\,d\rho\right)^{1/b'}
 \lesssim\big\|\|G_\rho\|_{L^a(S^{n-1})}\big\|_{L^{b'}([1,2])}.
\]
Since \(a\ge b'\) and \(\rho^{n-1}\) is comparable to \(1\) on this interval, the right side is bounded by a constant times \(\|G\|_{L^a(\Sigma)}\).
